\newcommand{\titulus}{\nomenFesti{Beatæ Mariæ Virginis Perdolentis.}
\dies{Die 15. Septembris.}}
\newcommand{\oratio}{\pars{Oratio.}

\noindent Deus, qui Fílio tuo in cruce exaltáto compatiéntem Matrem astáre voluísti, da Ecclésiæ tuæ, ut, Christi passiónis cum ipsa consors effécta, eiúsdem resurrectiónis párticeps esse mereátur.

\pars{Pro pace in universo mundo.} \scriptura{Sir. 50, 25; 2 Esdr. 4, 20; \textbf{H416}}

\vspace{-4mm}

\antiphona{II D}{temporalia/ant-dapacemdomine.gtex}

\vfill

\noindent Deus, a quo sancta desidéria, recta consília et iusta sunt ópera: da servis tuis illam, quam mundus dare non potest, pacem; ut et corda nostra mandátis tuis dédita, et hóstium subláta formídine, témpora sint tua protectióne tranquílla.

\noindent Per Dóminum nostrum Iesum Christum, Fílium tuum, qui tecum vivit et regnat in unitáte Spíritus Sancti, Deus, per ómnia sǽcula sæculórum.

\noindent \Rbardot{} Amen.}
\newcommand{\invitatorium}{\pars{Invitatorium.}

\vspace{-4mm}

\antiphona{IV}{temporalia/inv-salvatoremmundi.gtex}}
\newcommand{\hymnusmatutinum}{\pars{Hymnus.} \scriptura{Iacobus de Benedetti (\olddag{} 1306)}

\cuminitiali{II}{temporalia/hym-StabatMater.gtex}}
\newcommand{\matutinum}{\pars{Psalmus 1.} \scriptura{Lam. 1, 18; \textbf{H225}}

\vspace{-4mm}

\antiphona{VII d\textsuperscript{2}}{temporalia/ant-attenditeuniversi.gtex}

%\vspace{-2mm}

\scriptura{Ps. 23}

%\vspace{-2mm}

\initiumpsalmi{temporalia/ps23-initium-vii-d2-auto.gtex}

%\vspace{-1.5mm}

\input{temporalia/ps23-vii-d2.tex} \Abardot{}

\vfill
\pagebreak

\pars{Psalmus 2.} \scriptura{Ct. 1, 12}

\vspace{-4mm}

\antiphona{II D}{temporalia/ant-fasciculusmyrrhae.gtex}

\vspace{-2mm}

\scriptura{Ps. 45}

\vspace{-2mm}

\initiumpsalmi{temporalia/ps45-initium-ii-D-auto.gtex}

%\vspace{-1.5mm}

\input{temporalia/ps45-ii-D.tex} \Abardot{}

\vfill
\pagebreak

\pars{Psalmus 3.} \scriptura{Ct. 8, 7; \textbf{H300}}

\vspace{-4mm}

\antiphona{VIII C\textsuperscript{2}}{temporalia/ant-aquaemultae.gtex}

%\vspace{-2mm}

\scriptura{Ps. 86}

%\vspace{-2mm}

\initiumpsalmi{temporalia/ps86-initium-viii-C2-auto.gtex}

\input{temporalia/ps86-viii-C2.tex} \Abardot{}

\vfill
\pagebreak}
\newcommand{\matversus}{\noindent \Vbardot{} María conservábat ómnia verba hæc.

\noindent \Rbardot{} Cónferens in corde suo.}
\newcommand{\absolutio}{\cuminitiali{}{temporalia/absolutio-precibus.gtex}}
\newcommand{\benedictioi}{\cuminitiali{}{temporalia/benedictio-solemn-noscum.gtex}}
\newcommand{\benedictioii}{\cuminitiali{}{temporalia/benedictio-solemn-ipsavirgo.gtex}}
\newcommand{\benedictioiii}{\cuminitiali{}{temporalia/benedictio-solemn-pervirginem.gtex}}
\newcommand{\lectioi}{\pars{Lectio I.} \scriptura{Tob. 3, 7-15}

\noindent De libro Thobis.

\noindent Cóntigit Saræ fílias Ráguel, qui erat Ecbátanis Médiæ, ut ipsa audíret impropéria ab una ex ancíllis patris sui, quóniam trádita erat viris septem, et Asmódeus dæmónium nequíssimum occidébat eos, ántequam cum illa fíerent, sicut est sólitum muliéribus.

\noindent Et dixit illi ancílla: «Tu es, quæ suffócas viros tuos! Ecce iam trádita es viris septem, et némine eórum frúita es. Quid nos flagéllas causa virórum tuórum, quia mórtui sunt? Vade cum illis, nec ex te videámus fílium aut fíliam in perpétuum.»

\noindent In illa die contristáta est ánimo puélla et lacrimáta est et ascéndens in superiórem locum patris sui vóluit láqueo se suspéndere.

\noindent Et cogitávit íterum et dixit: «Ne forte impróperent patri meo et dicant: "Unicam habuísti fíliam caríssimam, et hæc láqueo se suspéndit ex malis"; et dedúcam senéctam patris mei cum tristítia ad ínferos. Utílius mihi est non me láqueo suspéndere, sed deprecári Dóminum ut móriar et iam impropéria non áudiam in vita mea. »

\noindent Porréctis mánibus ad fenéstram, Sara deprecáta est et dixit: «Benedíctus es, Dómine Deus miséricors, et benedíctum est nomen tuum sanctum et honorábile in sǽcula. Benedícant tibi ómnia ópera tua in ætérnum. Et nunc, Dómine, ad te fáciem meam et óculos meos diréxi.»

\noindent «Iube me dimítti désuper terram, et ne áudiam iam impropéria. Tu scis, Dómine, quóniam munda sum ab omni immundítia viri et non coinquinávi nomen meum neque nomen patris mei in terra captivitátis meæ. Unica sum patri meo, et non habet álium fílium qui possídeat hereditátem illíus, neque frater est illi próximus neque propínquus illi, ut custódiam me illi uxórem. Iam periérunt mihi septem, et ut quid mihi adhuc vívere? Et si non tibi vidétur, Dómine, occídere me, ímpera ut respiciátur in me et misereátur mei, et ne iam impropérium áudiam.»}
\newcommand{\responsoriumi}{\pars{Responsorium 1.} \scriptura{\Rbardot{} Tob. 3, 15 \Vbardot{} ibid. 3, 2; \textbf{H407}}

\vspace{-5mm}

\responsorium{I}{temporalia/resp-petodomine-CROCHU.gtex}{}}
\newcommand{\lectioii}{\pars{Lectio II.} \scriptura{Sermo in dom. infra oct. Assumptionis, 14-15: Opera omnia, Edit. Cisterc. 5 [1968], 273-274}

\noindent Ex Sermónibus sancti Bernárdi abbátis.

\noindent Martýrium Vírginis, tam in Simeónis prophetía quam in ipsa domínicæ passiónis história commendátur. Pósitus est hic, ait sanctus senex de párvulo Iesu, in signum cui contradicétur, et tuam ipsíus ánimam, ad Maríam autem dicébat, pertransíbit gládius.

\noindent Vere tuam, o Beáta mater, ánimam gládius pertransívit. Alióquin nónnisi eam pertránsiens, carnem Fílii penetráret. Et quidem posteáquam emísit spíritum tuus ille Iesus, ómnium quidem sed speciáliter tuus, ipsíus plane non áttigit ánimam crudélis láncea, quæ ipsíus, nec mórtuo parcens cui nocére non posset, apéruit latus, sed tuam útique ánimam pertransívit. Ipsíus nimírum ánima iam ibi non erat, sed tua plane nequíbat avélli. Tuam ergo pertransívit ánimam vis dolóris, ut plus quam mártyrem non immérito prædicémus, in qua nimírum corpóreæ sensum passiónis excésserit compassiónis efféctus.}
\newcommand{\responsoriumii}{\pars{Responsorium 2.} \scriptura{Lam. 1, 12; \textbf{H223}}

\vspace{-5mm}

\responsorium{VIII}{temporalia/resp-ovosomnes.gtex}{}}
\newcommand{\lectioiii}{\pars{Lectio III.}

\noindent An non tibi plus quam gládius fuit sermo ille, revéra pertránsiens ánimam, et pertíngens usque ad divisiónem ánimæ et spíritus: Múlier, ecce fílius tuus? O commutatiónem! Ioánnes tibi pro Iesu tráditur, servus pro Dómino, discípulus pro Magístro, fílius Zebedǽi pro Fílio Dei, homo purus pro Deo vero. Quómodo non tuam affectuosíssimam ánimam pertransíret hæc audítio, quando et nostra, licet sáxea, licet férrea péctora, sola recordátio scindit?

\noindent Non mirémini, fratres, quod María martyr in ánima fuísse dicátur. Mirétur qui non memínerit audísse se Paulum inter máxima géntium crímina memorántem quod sine affectióne fuíssent. Longe id fuit a Maríæ viscéribus, longe sit et a sérvulis suis.

\noindent Sed forte quis dicat: «Numquid non eum præscíerat moritúrum?». Et indubitánter. «Numquid non sperábat contínuo resurrectúrum?». Et fidénter. «Super hæc dóluit crucifíxum?». Et veheménter. Alióquin, quisnam tu, frater, aut unde tibi hæc sapiéntia, ut miréris plus Maríam compatiéntem quam Maríæ Fílium patiéntem? Ille étiam córpore mori pótuit, ista cómmori corde non pótuit? Fecit illud cáritas, qua maiórem nemo hábuit; fecit et hoc cáritas, cui post illam símilis áltera non fuit.}
\newcommand{\responsoriumiii}{\pars{Responsorium 3.} \scriptura{\Rbardot{} Iob 16, 17 \Vbardot{} Lam. 1, 12; \textbf{H219}}

\vspace{-5mm}

\responsorium{V}{temporalia/resp-caligaverunt-cumdox.gtex}{}}
\newcommand{\hymnuslaudes}{\pars{Hymnus.}

\cuminitiali{II}{temporalia/hym-EiaMater.gtex}}
\newcommand{\laudes}{\pars{Psalmus 1.} \scriptura{Ps. 62, 9}

\vspace{-4mm}

\antiphona{VIII G}{temporalia/ant-adhaesitanimamea.gtex}

%\vspace{-2mm}

\scriptura{Psalmus 62}

%\vspace{-2mm}

\initiumpsalmi{temporalia/ps62-initium-viii-G-auto.gtex}

%\vspace{-1.5mm}

\input{temporalia/ps62-viii-G.tex} \Abardot{}

\vfill
\pagebreak

\pars{Psalmus 2.} \scriptura{Phil. 1, 21; Gal. 6, 14; \textbf{H284}}

\vspace{-4mm}

\antiphona{I g}{temporalia/ant-mihivivere.gtex}

%\vspace{-2mm}

\scriptura{Canticum trium puerorum, Dan. 3, 57-88 et 56}

\initiumpsalmi{temporalia/dan3-initium-i-g-auto.gtex}

\input{temporalia/dan3-i-g-sinedox.tex}

\rubrica{Hic non dicitur Gloria Patri, neque Amen.}

\vfill

\antiphona{}{temporalia/ant-mihivivere.gtex}

\vfill
\pagebreak

\pars{Psalmus 3.}

\vspace{-4mm}

\antiphona{I d}{temporalia/ant-mortuussum.gtex}

%\vspace{-2mm}

\scriptura{Psalmus 149}

%\vspace{-2mm}

\initiumpsalmi{temporalia/ps149-initium-i-d-auto.gtex}

\input{temporalia/ps149-i-d.tex} \Abardot{}

\vfill
\pagebreak}
\newcommand{\lectiobrevis}{\pars{Lectio Brevis.} \scriptura{Col. 1, 24-25}

\noindent Nunc gáudeo in passiónibus pro vobis et adímpleo ea, quæ desunt passiónum Christi, in carne mea pro córpore eius, quod est ecclésia, cuius factus sum ego miníster secúndum dispensatiónem Dei, quæ data est mihi in vos, ut ímpleam verbum Dei.}
\newcommand{\responsoriumbreve}{\pars{Responsorium breve.}

\cuminitiali{VI}{temporalia/resp-pertesalutem.gtex}}
\newcommand{\benedictus}{\pars{Canticum Zachariæ.} \scriptura{Lc. 2, 35}

\vspace{-4mm}

\antiphona{VI F}{temporalia/ant-tuamipsiusanimam.gtex}

\vspace{-2mm}

\scriptura{Lc. 1, 68-79}

\vspace{-2mm}

\cantusSineNeumas
\initiumpsalmi{temporalia/benedictus-initium-vi-F-auto.gtex}

%\vspace{-1.5mm}

\input{temporalia/benedictus-vi-F.tex} \Abardot{}}
\newcommand{\benedicamuslaudes}{\cuminitiali{I}{temporalia/benedicamus-festis-bmv.gtex}}
\include{hebdomadaxxiv}
\include{feriaiii}
